\begin{frame}[t,label=B-why-pmes]{PME splits: one survives correction}
  \FrameTop
  % WTP basis: group-specific pay slopes (pme_channels.csv, block wtp / fiwtp); values in % of pay given up, group A / group B
  {\TableFont\setlength{\tabcolsep}{2.5pt}\renewcommand{\arraystretch}{1.05}%
  \begin{tabularx}{\linewidth}{@{}>{\RaggedRight}X>{\centering}p{1.6cm}>{\centering}p{1.3cm}>{\centering}p{1.3cm}>{\centering}p{1.3cm}>{\centering}p{1.3cm}>{\centering}p{1.0cm}>{\centering}p{1.0cm}>{\centering\arraybackslash}p{0.82cm}@{}}
    \toprule
    % [S7] attribute levels: 1-2 and 3-4 WFH days a week
    \textbf{Split (A / B)} & \textbf{n} & \textbf{1--2 days} & \textbf{3--4 days} & \textbf{Remote}\textsuperscript{a} & \textbf{Compr.} & \textbf{p} & \textbf{p compr.} & \textbf{q}\\
    \midrule
    Career penalty: y / n & 455 / 565 & 2 / 10 & 2 / 14 & 6 / 16 & 10 / 6 & .002 & .15 & .012\\
    Commute $>$ / $\le$45 min & 305 / 681 & 6 / 8 & 10 / 8 & 16 / 11 & 6 / 8 & .015 & .54 & .051\\
    WFH now: y / n & 666 / 338 & 9 / 5 & 12 / 5 & 14 / 9 & 8 / 6 & .094 & .50 & .22\\
    Pay $>$ / $\le$S\$6,000 & 408 / 435 & 8 / 7 & 11 / 10 & 15 / 10 & 9 / 6 & .77 & .33 & .77\\
    Public / local / MNC & 171/383/412 & 8/9/7 & 6/11/9 & 15/12/11 & 7/7/8 & .68 & .93 & .77\\
    Tenure: 3+ / $<$3 years & 805 / 199 & 7 / 6 & 9 / 12 & 12 / 11 & 7 / 8 & .62 & .69 & .77\\
    Mothers, child 0--6 / rest & 108 / 912 & 11 / 7 & 17 / 9 & 16 / 12 & 7 / 7 & .33 & 1.00 & .58\\
    \bottomrule
  \end{tabularx}\par}
  \vspace{0.08cm}
  \Body{\textbf{Reading:} only the career-concern split survives correction (location); about 40\% of its gap could be choice noise, and it was asked after the choices.}
  \SourceNote{\% of pay given up, group A / group B. p: test that the three location values are equal across groups (group-specific pay slopes); p compr.: same for the compressed week; q: Benjamini--Hochberg over the splits. Compr. = compressed week. \textsuperscript{a}Fully remote vs onsite with a moderate commute.}
  \BackToTwo{M-why-pmes}{Career concerns}{M-groups}{PME groups}
\end{frame}

\begin{frame}[t,label=B-emp-groups]{No group difference established}
  % [Q14, Q14b] equal-pay shares by group; [E18] p = .43 / .09 / .66; 198 contrasts, min q = .70; [Q14] p = .047 nominal
  % [M9] access row: compressed week 58.7 vs 58.1 at equal pay, p = .86; E19 not robust; [M7] further splits on B-why-pmes
  \FrameTop
  \Fig{F_emp_groups}{5.2cm}
  \SourceNote{Joint tests that all pay values are equal across groups: gender p = .43, age .09, care .66; none of 198 contrasts survives correction. \textsuperscript{1}Vs onsite with a moderate commute (group SEs provisional). Detail: \hyperlink{B-ra-gender}{\textcolor{Navy}{gender}}, \hyperlink{B-ra-age}{\textcolor{Navy}{age}}, \hyperlink{B-ra-care}{\textcolor{Navy}{care}}, \hyperlink{B-ra-role}{\textcolor{Navy}{role}}, \hyperlink{B-ra-edu}{\textcolor{Navy}{edu}}.}
  \BackToTwo{M-groups}{PME groups}{M-why-pmes}{Career concerns}
\end{frame}

\begin{frame}[t,label=B-access]{Valued more where offered? Not robust}
  % [E19] status X, NOT ROBUST (Lead decision 5); [M9] 58.7 vs 58.1 at equal pay, p = .86; +5.5 [-4.4, 15.4], p = .27;
  % slope 2.79 vs 4.70; [M20] cheaper option in 83.2% of PME pairs; [D2, D2x] take-up
  \FrameTop
  {\centering\fontsize{8.5}{10}\selectfont\input{tables/T_access}\par}
  \SourceNote{The gap is statistically significant only under a common pay slope. PMEs with access are less pay-sensitive (log-pay coefficient 2.79 vs 4.70); the compressed week was cheaper in 83\% of PME pairs. Experience and sorting cannot be told apart.}
  \BackToTwo{M-puzzle}{Puzzle}{M-groups}{Groups}
\end{frame}

\begin{frame}[t,label=B-emp-robust]{PME values stay positive}
  \FrameTop
  \Tab{T_emp_robust}
  % [E15, RB2] moves of about 1 point (same-side part-time p = .06); [E8] job security X
  \SourceNote{Outliers: implausible salaries dropped; same-side: always chose one side. Values stay positive (staggered \ap 0) and significant at 5\%, except part-time without \hyperlink{B-ra-sameside}{\textcolor{Navy}{same-side}} respondents; most move \ap 1 point (\hyperlink{B-ra-trim}{\textcolor{Navy}{trimming}}).}
  \BackToTwo{M-robust}{Robustness}{M-pmes}{PMEs}
\end{frame}

\begin{frame}[t,label=B-org-robust]{Employer values hold in every sample}
  \FrameTop
  % appended_wtp.xlsx sheets main, robust_outlier (salary_out excluded, N = 392), robust_trim (top and bottom 1% of income, N = 391), robust_same_side (N = 381); pay cut = minus WTP
  \Tab{T_org_robust}
  \SourceNote{Cuts for 1--2 days, 3--4 days, fully remote and reduced load or part-time stay statistically significant in every sample; the compressed week stays near 7 and is not distinguishable from zero. Moves are \ap 1 point, up to 2.5 without same-side respondents. Main results use all 397 firms.}
  \BackToTwo{M-robust}{Robustness}{M-employers}{Employers}
\end{frame}

% ==========================================================================
% Employers
% ==========================================================================
% [R11, R11b, Q8] high turnover risk 34% (appended); q over 7 comparisons
\BackupTab{B-org-salary}{Why employer \% is secondary}{T_org_salary}{M-employers}{Employers}{Employers' pay-cut sizes rest on a weakly identified pay slope; the ranking is firmer.}

\begin{frame}[t,label=B-why-org]{Employer reservations: two preliminary contrasts}
  \FrameTop
  \setbeamerfont{itemize/enumerate body}{size=\fontsize{9.5}{11.6}\selectfont}
  % polish pass (BKb): the main-slide table style (MainTableStyle, 10 pt, for projection); panel rows bold
  % Ink (Navy is structure only); the last column fits the bold Compressed header
  % final fixer (7 Oct 2026): notes print at 8.0 pt (were 6.9 pt); the table's arraystretch is 1.08 here (\MainTableStyle: 1.15)
  % Appended sample: pay cut required (= minus WTP) from appended_wtp.xlsx (physical_presence, concern_creativity) / data/wtp_deck.csv;
  % p: ratio test, location WTPs equal across the two groups (group-specific pay slopes); q: Benjamini--Hochberg over 11 split tests
  {\MainTableStyle\renewcommand{\arraystretch}{1.08}%
  \begin{tabularx}{\linewidth}{@{}>{\RaggedRight}X>{\centering}p{1.1cm}>{\centering}p{1.7cm}>{\centering}p{1.7cm}>{\centering}p{1.7cm}>{\centering\arraybackslash}p{2.35cm}@{}}
    \toprule
    % [S7] attribute levels: 1-2 and 3-4 WFH days a week
    \textbf{Pay cut required, \% of pay} & \textbf{Firms} & \textbf{1--2 days} & \textbf{3--4 days} & \textbf{Remote} & \textbf{Compressed}\\
    \midrule
    \multicolumn{6}{@{}l}{\textbf{Do most PME roles need on-site presence?}\enspace p = .025,\ q = .14}\\
    \quad Yes, most roles & 137 & \textbf{\color{Employer}52} & 60 & 67 & 1\\
    \quad Only some or few & 260 & \textbf{\color{Employer}3} & 25 & 40 & 8\\
    \midrule
    \multicolumn{6}{@{}l}{\textbf{Worried that creativity suffers when people are apart?}\enspace p = .14,\ q = .43}\\
    \quad Yes & 201 & 24 & 44 & \textbf{\color{Employer}61} & 18\\
    \quad No & 196 & 15 & 28 & \textbf{\color{Employer}36} & $-$3\\
    \bottomrule
  \end{tabularx}\par}
  \vspace{0.06cm}
  \begin{itemize}\setlength{\itemsep}{0.04cm}
    \item \textbf{The presence contrast is significant at the 5\% level} (p = .025) but not after correcting for the 11 split tests (q = .14); the creativity contrast is not (p = .14).
    % EB s9 l.182: presence and concern items sit in Block 5, after the DCE block
    \item \textbf{Asked after the choices:} the answers may partly rationalise the hiring choices.
    \item \textbf{Separate pay slope for each group:} values are group-specific, so they are preliminary associations.
  \end{itemize}
  \SourceNote{Pay cut employers would require of a candidate wanting the arrangement, \% of pay (minus the model's WTP). p: test that the three location values are equal across the two groups; q: Benjamini--Hochberg over 11 split tests. Many values go beyond the pay gaps shown (17--33\%): read the contrasts, not the sizes. \hyperlink{B-why-org-tests}{\textcolor{Navy}{All split tests}}.}
  \BackTo{M-where}{Reservations}
\end{frame}

% [Q10-Q13] pay cut required by employer group (wtp_deck.csv, org splits); p ratio test, q = BH over 11 split tests
% EB s10 QID22_6: a survey item, labelled salary offset, never FWAs let firms pay less (fixer pass 7 Oct 2026, worker-03)
\BackupFig{B-org-splits}{All employer splits}{F_org_splits}{5.4cm}{M-where}{Reservations}{Pay cut required, 95\% CI. Presence and no WFH today: p $\approx$ .02 each, q = .14 over all \hyperlink{B-why-org-tests}{\textcolor{Navy}{11 split tests}}; creativity p = .14, q = .43. Preliminary: asked after the choices. Salary offset: agrees that offering FWAs reduces the salary needed to attract and retain PMEs (survey item).}

% [B2] agreement shares; [Q11, R14b, M12] creativity p = .003, q about .06; Lead decision 6: concerns, not beliefs
% EB s9 l.183: no free-slope check yet for the presence and creativity splits; EB s9 l.182: items asked after the choices
\BackupTab{B-org-concerns}{Concerns are common; creativity tracks choices}{T_concerns}{M-where}{Reservations}{Agree = 4--5 of 5. p: joint WFH test, one pay slope. Preliminary associations: items asked after the choices. Perceived vs true cost: not separable.}

\begin{frame}[t,label=B-org-view]{Employers: talent upside, practical hurdles}
  % [B1, B2, B3]; [B1] fairness top challenge at 200 or fewer staff, 28% (shown in the figure); [P16] MOM PQ 2937; [S3] 12%; [R15, Q13, M18] no size association
  \FrameTop
  \Fig{F_org_view}{5.2cm}
  \SourceNote{We detect no size association in hiring choices; compressed-week offering rises with firm size (14, 18, 26\%; p = .07). MOM (PQ, 10 Sep 2026): in 2025, nearly 7 in 10 SMEs and about 8 in 10 non-SMEs offered a scheduled FWA; definitions differ, so not comparable with the 66.4\% series.}
  \BackToTwo{M-employers}{Employers}{M-where}{Reservations}
\end{frame}

\begin{frame}[t,label=B-industry-note]{Why we show no industry estimates}
  % EB s10: cells of 1-44 respondents, 13 of 17 industries below 25; [R15] four largest industries p = .87
  % [M10, M12] presence and creativity splits; EB s9 l.183: free-slope check pending
  \FrameTop
  \begin{itemize}\setlength{\itemsep}{0.12cm}
    \item Industry cells hold 4 to 60 employers (11 of 17 below 25); some intervals are wider than the scale.
    \item Four largest industries (ICT, finance, manufacturing, professional services): we detect no difference in WFH valuations (p = .94).
    \item Reservations are associated with role needs and concerns (presence; creativity), as preliminary associations under a common pay-slope model; industry differences cannot be ruled out.
  \end{itemize}
  \SourceNote{Industry mix: \hyperlink{B-industry}{\textcolor{Navy}{Industries}}.}
  \BackTo{M-where}{Reservations}
\end{frame}

% ==========================================================================
% Where the sides meet
% ==========================================================================
\begin{frame}[t,label=B-modelfree]{Raw choices, before any model}
  % [E13, E14, R13, M21]
  \FrameTop
  \Tab{T_modelfree}
  \SourceNote{Raw choice shares; the other attribute varies and cells are fixed by version: read as direction.}
  \BackToTwo{M-meet}{Two sides}{M-employers}{Employers}
\end{frame}

\begin{frame}[t,label=B-puzzle]{Compressed week: design balance, offering, use}
  \FrameTop
  % polish pass (BKb): text lines at 9 pt (dense two-column backup), bold table headers, wider right
  % table so the size row value stays on one line
  \begin{columns}[T,onlytextwidth]
    \begin{column}{0.41\textwidth}
      \PanelTitle{Navy}{Design: usually the cheaper option}
      % [M20] cheaper option in 83.2% of PME pairs and 90.8% of employer pairs
      {\fontsize{9}{11}\selectfont\RaggedRight The lower-paid option in 83\% of PME and 91\% of employer pairs.\par}
      \vspace{0.06cm}
      {\TableFont\setlength{\tabcolsep}{3pt}%
      \begin{tabular}{@{}>{\RaggedRight}p{3.15cm}>{\centering}p{1.15cm}>{\centering\arraybackslash}p{1.15cm}@{}}
        \toprule
        \textbf{\% chosen if} & \textbf{cheaper} & \textbf{dearer}\\
        \midrule
        % [M21] PMEs: compressed-week job 27.5% when cheaper, 85.5% when dearer
        PMEs: compressed & 27.5 & 85.5\\
        % [M21] PMEs: other options 24.9% and 75.1%
        PMEs: others & 24.9 & 75.1\\
        % [M21] employers: compressed-week candidate 66% cheaper, 36% dearer with n = 33
        Firms: compressed & 67 & 37\textsuperscript{c}\\
        % [M21] employers: other candidates 55% and 45%
        Firms: others & 56 & 44\\
        \bottomrule
      \end{tabular}\par}
      \vspace{0.06cm}
      % [M21] PMEs: positive value in both price positions; employers mixed; [M20] employer no-penalty identified mainly from cheaper candidates
      {\fontsize{9}{11}\selectfont\RaggedRight \textbf{PMEs:} consistent with a positive value at either price. \textbf{Firms:} mixed; ``no detectable penalty'' rests mainly on cheaper candidates.\par}
    \end{column}
    \begin{column}{0.56\textwidth}
      \PanelTitle{Navy}{Surveyed firms offering it (70 of 397)}
      {\TableFont\setlength{\tabcolsep}{3pt}%
      \begin{tabular}{@{}>{\RaggedRight}p{2.85cm}>{\centering}p{2.55cm}>{\centering\arraybackslash}p{2.35cm}@{}}
        \toprule
        \textbf{Trait} & \textbf{\% offer} & \textbf{OR [95\% CI]}\\
        \midrule
        % [M17] 26.5% (53/200) vs 9.7% (10/103); OR 2.94 [1.37, 6.33], p = .006
        Written policy & 21.4 vs 9.1 & \textbf{2.41} [1.20, 4.87]\\
        % [M17] offering part-time OR 3.19 [1.65, 6.18], p < .001
        Offers part-time\textsuperscript{a} & & \textbf{3.03} [1.67, 5.51]\\
        % [M18] most roles need presence OR 1.71 [0.94, 3.12], p = .080
        Presence need & & 1.53 [0.88, 2.66]\\
        % [M18] fully tech-ready OR 1.20 [0.66, 2.19]
        Fully tech-ready & & 0.83 [0.48, 1.43]\\
        % [M18] productivity worry OR 1.36 [0.59, 3.10]
        Productivity concern\textsuperscript{b} & & 1.61 [0.77, 3.37]\\
        % [M18] hiring difficulty 21.5% vs 20.0%, p = .74 (CSV chi2_p .7448)
        Hard to hire & 18.4 vs 16.8 & p = .69\\
        % [M18] by size band 18.9 / 19.8 / 25.0, p = .57
        Size: $\le$200 / 201--1,000 / 1,001+ & 13.7 / 17.9 / 25.6 & p = .066\\
        \bottomrule
      \end{tabular}\par}
      \vspace{0.06cm}
      {\fontsize{9}{11}\selectfont\RaggedRight \textbf{Surveyed PMEs:} access 21\%; take-up 31\%; career concern OR 0.86 [0.44, 1.67].\par}
    \end{column}
  \end{columns}
  \vspace{0.08cm}
  % EB s9: hiring choices do not measure operating costs; COMMENTS B-puzzle row: pilot protocol on B-pilot
  {\fontsize{9}{11}\selectfont\RaggedRight Operating costs unmeasured: \hyperlink{B-pilot}{\textcolor{Navy}{pilot checklist}}.\par}
  % [M20] only 33 dearer employer pairs
  \SourceNote{Logit on size, WFH now, policy, tech, presence (N = 397); ORs are not risk ratios; policy may follow offering. \textsuperscript{a}With other offerings (policy OR 2.16 [1.04, 4.49]). \textsuperscript{b}With concerns, talent view, hiring, turnover: ORs 0.75--1.61, p $\ge$ .16. \textsuperscript{c}41 pairs.}
  \BackToTwo{M-puzzle}{Compressed week}{M-discussion}{Zones}
\end{frame}

\begin{frame}[t,label=B-pilot]{A compressed-week pilot: checklist}
  % Fixer pass 7 Oct 2026, merging verified persona fixes: worker-02 (pay and leave unchanged; collective
  % agreements; voluntary for each employee, reversible without penalty; care hours; appraisal by output;
  % worker data only with consent; stop criteria agreed with the partners), employer-06 (for firms and staff;
  % fairness to staff who cannot take part; operating costs and coverage as outcomes), editor-06 (comparison
  % group), mom_ps-07 (bounds to agree), worker-07 (eligibility and fairness). Items at 9 pt and columns 0.34/0.24/0.39 (polish pass: 0.32/0.27/0.37):
  % at 10 pt (worker-02's fallback) the frame still overflowed by 1.28 cm in layout_sim; this is a backup frame.
  % [S7] compressed week = same total hours over fewer, longer days; never a shorter week
  % [B1] fairness to staff who cannot have FWAs: 23%, the top challenge at 200 or fewer staff (28%); [B2] 50% agree
  % it is unfair to offer FWAs to PMEs when frontline staff cannot have them (not printed)
  \FrameTop
  \setbeamerfont{itemize/enumerate body}{size=\fontsize{9}{11}\selectfont}
  \begin{columns}[T,onlytextwidth]
    \begin{column}{0.32\textwidth}
      \PanelTitle{Navy}{Design}
      \begin{itemize}\setlength{\itemsep}{0.04cm}
        \item Same hours, four longer days; pay and leave unchanged
        \item Legal, contract and collective-agreement review
        \item Voluntary for firms and staff; reversible without penalty
        \item Eligibility; fairness to staff who cannot join
      \end{itemize}
    \end{column}
    \begin{column}{0.27\textwidth}
      \PanelTitle{Navy}{Safeguards}
      \begin{itemize}\setlength{\itemsep}{0.04cm}
        \item Service and coverage plan
        \item Fatigue, safety and care hours
        \item Workload spillovers
        \item Appraisal by output, not presence
      \end{itemize}
    \end{column}
    \begin{column}{0.37\textwidth}
      \PanelTitle{Navy}{Evaluation}
      \begin{itemize}\setlength{\itemsep}{0.04cm}
        \item Outcomes: output, operating costs, coverage, retention, take-up, worker experience
        \item Stop and scale criteria agreed with the partners in advance
        \item Independent evaluation with a comparison group; worker data only with consent, never for appraisal
        % Governance: publication rule not yet agreed (mandate; EB G0)
        \item Publication rule: [PI to confirm]
      \end{itemize}
    \end{column}
  \end{columns}
  % polish pass (BKb): a hairline Rule sets the summary lines apart from the checklist; lines at the
  % 9.5 pt backup body (Body macro)
  \vspace{0.10cm}{\color{Rule}\hrule height 0.4pt}\vspace{0.16cm}
  % [Q5] PMEs value it, employers show no detectable hiring penalty; [M19] no firm group penalises it;
  % [M20] it was usually the cheaper option; mandate wording for the compressed-week conclusion
  \Body{\textbf{Why a pilot:} most promising common-ground test; evidence moderate; operating costs unmeasured.}
  \vspace{0.06cm}
  % mom_ps-07: the bounds are for the partners to agree, as on B-optC; no governance fact is implied
  \Body{\textbf{To agree:} number of firms, duration, comparison group, funder, data access, report-back date.}
  \vspace{0.06cm}
  % M-options, Option A owners: SNEF with volunteer firms; NTUC co-designs safeguards; TAFEP guidance
  \Body{\textbf{Possible leads:} SNEF with volunteer firms; NTUC co-designs safeguards; TAFEP guidance.}
  % [S7] survey wording, both questionnaires: same total hours across fewer days
  % EB s9: the employer task is a hire, not an operating-cost measure
  \SourceNote{Both surveys: the same total hours over fewer, longer days. The hiring choices did not record rostering, coverage or fatigue.}
  \BackToTwo{M-puzzle}{Compressed week}{M-options}{Options}
\end{frame}

% Round 5: the table prints employer % discounts [R6, R7, G6], so the frame carries the model-dependent header
% and should not be linked from a short-path main slide (ROUND5_SPEC s5); B-optB links here.
\begin{frame}[t,label=B-flexiload]{Flexi-load: mostly pay for hours}
  \FrameTop % first body token (centring; never a frame subtitle)
  {\centering\Caveat{Model-dependent; not a pay recommendation and not a main finding.}}
  \Tab{T_flexiload}
  % [S7] 4 of 5 days; [R6, R7, G6] about pro-rata; [M13] part-time 35% where neither concern; [P11] PTRG under review, no enhancement by 6 Oct 2026
  % Declutter: [P5] 5.9% vs 13.7% part-time share of all women 25-64 is on B-laboursupply.
  \SourceNote{4 of 5 days = 20\% fewer hours, so employer discounts near pro-rata show no sign of an extra flexibility penalty. Part-time draws a larger required cut (about 36\%) because it cuts more hours. MOM is reviewing the PTRG for flexi-load (COS 2026); \hyperlink{B-laboursupply}{\textcolor{Navy}{part-time context}}.}
  \BackToTwo{M-options}{Options}{M-where}{Reservations}
\end{frame}

% ==========================================================================
% Barriers
% ==========================================================================
\begin{frame}[t,label=B-career-who]{Who expects a career penalty?}
  % F_career_who is drawn wide (14.0 x 4.6 cm), so it runs full width at scale 1.0.
  % It shows: share agreeing that shorter or more flexible hours would hurt their career, by perceived job-loss band
  % [M24] 28-33% to 50-62%; by familiarity with the Guidelines [M25] 45-54% vs 22%, odds ratio about halved
  % full approval p = .44; all-PME dotted band 37-45% [M23]
  \FrameTop
  \Fig{F_career_who}{4.6cm}
  \vspace{0.12cm}
  % [P15] MOM PQ 2936, 9 Sep 2026: about 6% of resident FWA users reported setbacks in 2025
  {\fontsize{9}{11}\selectfont\RaggedRight \textbf{Parents of 0--6s:} +12 pp after controls (p = .001). \textbf{MOM:} about 6\% of resident FWA users reported a setback in 2025 (reported experience, not expectation).\par}
  \SourceNote{Agree (4--5 of 5) that shorter or more flexible hours would hurt their career: 45\% of surveyed PMEs. Job loss: self-rated 12-month chance. MOM PQ 2936 (9 Sep 2026); not comparable. \hyperlink{B-career}{\textcolor{Navy}{Chart}}.}
  \BackToTwo{M-barriers}{Barriers}{M-why-pmes}{Career concerns}
\end{frame}

\begin{frame}[t,label=B-career]{Over a third of surveyed PMEs expect a career penalty}
  % [B5, M23] 45 [37]; S4 panel composition of parents of 0-6; [P8] figure note: workers in general, not PMEs (FACTCHECK F22, F30)
  \FrameTop
  \Fig{F_career}{4.7cm}
  \SourceNote{Agree = 4--5 of 5: an expected penalty, not a reported setback. PAP--NTUC survey: all workers, not PMEs. \hyperlink{B-career-who}{\textcolor{Navy}{Who expects it}}.}
  \BackToTwo{M-barriers}{Barriers}{M-why-pmes}{Why PMEs}
\end{frame}

\begin{frame}[t,label=B-tgfwar]{What we can say about the Guidelines}
  \FrameTop
  \setbeamerfont{itemize/enumerate body}{size=\fontsize{9.5}{11.6}\selectfont}
  \begin{columns}[T,onlytextwidth]
    \begin{column}{0.48\textwidth}
      \PanelTitle{Navy}{What we observe (2026)}
      % [D6, D6x, D7, D7b, D5x, D10]
      \begin{itemize}\setlength{\itemsep}{0.08cm}
        \item Familiar: 55\% of surveyed PMEs, 29\% of juniors, 58\% of employers
        \item Of PMEs reporting a formal request, 4\% were declined (7\%, stricter definition): a lower bound
        \item More FWAs came via company policy (43\%) than individual request (30\%)
        \item Aware firms with a written policy got more requests (68\% vs 19\%; an association)
      \end{itemize}
    \end{column}
    \begin{column}{0.48\textwidth}
      \PanelTitle{Navy}{What we cannot say}
      \begin{itemize}\setlength{\itemsep}{0.08cm}
        \item Whether the Guidelines changed outcomes (no before-and-after)
        \item National rates (unweighted panels)
        \item Why PMEs do not ask (not asked)
        \item How firms decide staff requests (experiment: hiring)
      \end{itemize}
    \end{column}
  \end{columns}
  \SourceNote{See \hyperlink{B-approvals}{\textcolor{Navy}{two views}}, \hyperlink{B-approvals-tab}{\textcolor{Navy}{table}}, \hyperlink{B-awareness}{\textcolor{Navy}{awareness}}.}
  \BackTo{M-barriers}{Barriers}
\end{frame}

% [D6, D6x]; [P8] PAP Women's Wing-NTUC WAF survey: 58% aware, 29% know where to turn (FACTCHECK F30, V-21)
\BackupFig{B-awareness}{Awareness is wide but shallow}{F_awareness}{5.0cm}{M-barriers}{Barriers}{PAP Women's Wing--NTUC survey (online, Nov--Dec 2025; sampling not published): 58\% aware; 29\% know where to turn if a request is denied or mishandled.}

% [D7, D7c, D8] figure; [P3] about 94%, data for 2024, almost entirely before the Guidelines, undefined (FACTCHECK F11, F13)
% [P15] MOM PQ 2936, 9 Sep 2026: about 6% of resident FWA users reported a setback in 2025; experience, not expectation (moved off M-barriers, round 5)
\begin{frame}[t,label=B-approvals]{Same process, two vantage points}
  \FrameTop
  \Fig{F_approvals_two}{4.6cm}
  \vspace{0.12cm}
  % [P3] approval about 94%, data for 2024; [P15] about 6%, 2025 data; [B5, M23] expected penalty is the PMEs' expectation, not comparable
  {\fontsize{9}{11}\selectfont\RaggedRight \textbf{MOM, for comparison (definitions differ):} \ap 94\% approved (2024 data, almost entirely before the Guidelines). After use, \ap 6\% of resident FWA users reported a career setback (2025): experience, not PMEs' expectation.\par}
  \SourceNote{Different questions: PMEs report their own request; employers, their firm's approval share. Gaps may reflect partial approvals, firm vs individual rates, or managers' partial view.}
  \BackTo{M-barriers}{Barriers}
\end{frame}

% [D7, D7b, D9]
\BackupTab{B-approvals-tab}{Request outcomes in detail}{T_approvals}{M-barriers}{Barriers}{Many PMEs with automatic FWAs reported a formal request (\hyperlink{B-requestrate}{\textcolor{Navy}{why}}); both sides cover only requests made. MOM comparators: \hyperlink{B-approvals}{\textcolor{Navy}{two views}}.}

% [D9]; [P3] 18.5%, data for 2024
\BackupTab{B-requestrate}{Why no headline request rate}{T_requestrate}{M-barriers}{Barriers}{The item left ``formal'' undefined, with no ``automatic'' option. MOM (2024 data): 18.5\% of employees who needed an FWA made a formal request.}

% ==========================================================================
% Limits and next steps
% ==========================================================================
% One table: limit, what it means, next step (round 5: 60 design cells and the
% pending employer free-slope check added; COMMENTS slide 21).
\begin{frame}[t,label=B-limits]{Limits in detail}
  \FrameTop % first body token (centring; never a frame subtitle)
  % polish pass (BKb): one backup table style: TableFont, 4 pt column gaps, row spacing via the RowSep macro
  {\TableFont
  \setlength{\tabcolsep}{4pt}
  \renewcommand{\arraystretch}{1.10}
  \begin{tabularx}{\linewidth}{@{}>{\raggedright\arraybackslash}p{2.6cm}>{\raggedright\arraybackslash}X>{\raggedright\arraybackslash}p{3.5cm}@{}}
    \toprule
    \textbf{Limit} & \textbf{What it means} & \textbf{Next step}\\
    \midrule
    % [S6, S6b] every gap 17-33%; [R13] cheaper candidate chosen 57%; EB s3 header: employer salary term weakly identified
    \textbf{Pay gaps 17--33\%} & \% of pay is model-based; employers weakly pay-sensitive (cheaper won 58\%) & A small-gap wave\\ \RowSep
    % EB s9: 10 versions x 6 tasks = 60 distinct choice situations per survey; respondent-clustered SEs overstate precision
    \textbf{60 design cells} & 10 versions $\times$ 6 tasks per side: intervals likely too narrow & Cluster by version\\ \RowSep
    % [M20] 83.2% of PME and 90.8% of employer pairs
    \textbf{Compressed week} & Cheaper in 83\% (PME) and 91\% (employer) of its pairs & Balance pay positions\\ \RowSep
    % EB s9 l.183: no free-slope check yet; [M8, M19] free slopes only for the compressed week; EB s9 l.182 items asked after the choices
    \textbf{Employer splits} & Preliminary associations, asked after the choices & Ask these items before the choices\\ \RowSep
    % [S2, S5] a hire between two equally qualified candidates; Lead decision 7; EB s9: hire, not request
    \textbf{Employer task} & A hire, not a staff request; not an operating cost & Pilots that record operations\\ \RowSep
    \textbf{Samples} & Unweighted panels, mostly large firms & Items in MOM surveys\\ \RowSep
    % [D9] misread; [D7b] askers only, a lower bound; [B6] no why-not item
    \textbf{Request item} & Misread; refusals among askers only; no why-not item & Option D items\\
    \bottomrule
  \end{tabularx}\par}
  % [S9] respondents knew SNEF would use the results; EB s9: 118 of 1,020 = 12% capped; task-1 coding; [E15] choice results hold without the weaker panel
  \SourceNote{Also: stated choices (SNEF use known); no before-and-after; 12\% of salaries capped; task-1 coding. Choice results hold without outliers or same-side respondents. \hyperlink{B-nextsteps}{\textcolor{Navy}{Next steps}}.}
  \BackToTwo{M-limits}{Limits}{M-robust}{Robustness}
\end{frame}

\begin{frame}[t,label=B-nextsteps]{Sharpening the evidence}
  \FrameTop
  \setbeamerfont{itemize/enumerate body}{size=\fontsize{9.5}{11.6}\selectfont}
  \begin{columns}[T,onlytextwidth]
    \begin{column}{0.48\textwidth}
      \PanelTitle{Navy}{Before publication}
      \begin{itemize}\setlength{\itemsep}{0.08cm}
        % EB s9 l.183: employer free-slope equal-pay test for the presence and creativity splits, not yet run; [M8]
        \item Employer role and concern items asked before the choices
        \item Add a position term; fix early task-1 coding
        \item Cluster by survey version; wild bootstrap
        \item Log vs rank salary models
        % [B6] 68% of PMEs, 61% of employers wrote a comment
        \item Code open-text comments (68\% of PMEs, 58\% of employers wrote one)
      \end{itemize}
    \end{column}
    \begin{column}{0.48\textwidth}
      \PanelTitle{Navy}{A follow-up wave, if supported}
      \begin{itemize}\setlength{\itemsep}{0.08cm}
        % design proposal, not data (internal_drafts: B-nextsteps OK)
        \item Small pay gaps (0--15\%) for precise prices
        % [M20] balance the compressed week across pay positions
        \item Balance the compressed week across pay positions
        \item Ask why people do not request; reword the request item, with partial approval
        \item Information module (trial evidence; Guidelines; career outcomes)
        % [F1] 490 of 1,020 = 48% agreed to one short follow-up survey in 6-12 months; consent scope: one survey only; [F2] no employer re-contact item
        \item One short survey of the 490 PMEs (48\%) who agreed; more needs fresh consent (\hyperlink{B-followup}{\textcolor{Navy}{follow-up}})
      \end{itemize}
    \end{column}
  \end{columns}
  \BackToTwo{M-limits}{Limits}{M-options}{Options}
\end{frame}

\begin{frame}[t,label=B-followup]{Follow-up design, if supported}
  \FrameTop % first body token (centring; never a frame subtitle)
  % polish pass (BKb): backup table style (TableFont, 4 pt gaps, RowSep)
  % final fixer (7 Oct 2026): notes print at 8.0 pt (were 6.9 pt); rows at 1.0 and 1.5pt apart (were 1.05, 3pt) keep the button row
  {\TableFont\setlength{\tabcolsep}{4pt}\renewcommand{\arraystretch}{1.0}\renewcommand{\RowSep}{\addlinespace[1.5pt]}%
  \begin{tabularx}{\linewidth}{@{}>{\RaggedRight}p{1.85cm}>{\RaggedRight}X>{\RaggedRight\arraybackslash}X@{}}
    \toprule
     & \textbf{Aim 1: do stated choices predict behaviour?} & \textbf{Aim 2: what changes in firms that adopt?}\\
    \midrule
    \textbf{Outcomes} & FWA take-up, job moves, promotion, appraisal vs each PME's choices & Productivity, retention, hires, exits in volunteer pilot firms (Option A or C), vs a comparison group\\ \RowSep
    % [F1] 490 of 1,020 = 48%; about 5-minute survey in 6-12 months; window Feb-Aug 2027; [F2] employer survey asked no re-contact question
    \textbf{Consent now} & 490 of 1,020 PMEs (48\%) agreed to one short (about 5-minute) survey in 6--12 months (Feb--Aug 2027) & None: the employer survey asked no re-contact question\\ \RowSep
    \textbf{Partner roles} & \textbf{NTUC:} reach members beyond online panels; \textbf{MOM:} linkage only with fresh consent, under its data rules & \textbf{SNEF:} recruit fresh volunteer member firms; \textbf{MOM:} firm linkage only with consent\\ \RowSep
    % [V1] 56% matched up to three years later; [V6] quits 7.20% to 4.80%; [V8] six waves 2012-2024
    \textbf{Templates} & Maestas et al.: 56\% matched up to 3 years later (\hyperlink{B-validity}{\textcolor{Navy}{validity}}) & Trip.com hybrid trial: quits fell 7.20\% to 4.80\%; German LPP: consented employer--employee linkage, six waves 2012--2024\\
    \bottomrule
  \end{tabularx}\par}
  \vspace{0.08cm}
  % [F1] scope: one short follow-up survey only; not general re-contact, linkage or a firm study (mandate wording)
  {\fontsize{9}{11}\selectfont\RaggedRight
    \textbf{Consent scope:} 490 PMEs consented to one short survey in 6--12 months; further waves, linkage and firm-side work require fresh consent.\par % [F1] 490; one survey in 6-12 months
    \vspace{0.06cm}
    % Governance facts must come from the PI (mandate; EB G0): placeholders only
    \textbf{Data owner, funder and publication:} [PI to confirm] (\hyperlink{B-governance}{\textcolor{Navy}{governance}}).\par}
  \SourceNote{Consenters lean to caregivers (52.6\% vs 39.7\% agreed).}
  \BackTo{M-conclusion}{Conclusion}
\end{frame}

% ==========================================================================
% Options
% ==========================================================================
\OptionFrame{B-optA}{Option A: put compressed weeks on the menu and test them}
  {\textbf{Zone:} A, common ground\quad\textbf{Strength:} moderate; operating costs unmeasured}
  {
   % [M16] 20.8% offer; [M17] 26.5 vs 9.7, an association; [M20] 83.2 and 90.8
   PMEs would give up about 7\% of pay for it; employers would require about 7\% (wide interval): no detectable gap. 1 surveyed firm in 6 offers it, more often with a written policy (21\% vs 9\%). It was the cheaper option in 83\% (PME) and 91\% (employer) of its pairs.}
  {% W5 M-options owner: SNEF with volunteer firms; NTUC co-designs safeguards; TAFEP guidance
   \textbf{SNEF}, with volunteer member firms: put it on the written menu and test it in a bounded, evaluated pilot (\hyperlink{B-pilot}{\textcolor{Navy}{checklist}}). \textbf{NTUC:} co-design safeguards. \textbf{TAFEP:} guidance.}
  {% [P21] TAFEP, IHRP and SNEF resources (V-09)
   % TODO(PI): confirm the SNEF guide's exact title with SNEF (FACTCHECK V-09)
   TAFEP and IHRP resources; SNEF guidance on FWA requests.}
  {% [S7] the same total hours over four longer days
   Not shorter hours (same hours, four longer days), a mandate or a case for lower pay; not shown to be costless.}
  {% [D2x] take-up 31 [33]
   Output, coverage, rostering, fatigue; take-up where available (31\%).}
  {\SourceNote{Associations, not causes: a policy may follow offering. Hiring choices do not measure operating costs.}%
   \BackToTwo{M-options}{Options}{M-puzzle}{Puzzle}}

\OptionFrame{B-optB}{Option B: fit arrangements to the role, with flexi-load}
  {\textbf{Zone:} B, depends on the role\quad\textbf{Strength:} moderate; preliminary associations}
  {% [Q10, M10] 27 vs 48 (appended); [Q10b] 50 (appended); [Q10c] p = .002 vs 1-2 days; [M12] p = .006, q = .058 over 20 splits; EB s9 l.183 one pay slope, free-slope check pending
   % [M13] part-time 35 where neither presence need nor creativity concern
   Where most roles need presence, firms would require a pay cut of about 52\% for 1--2 WFH days (3\% elsewhere), against about 1\% for a compressed week. Presence: p = .025 (q = .14 over 11 split tests; preliminary).}
  {% [P11] Job Redesign+ up to 70% of qualifying costs for SMEs, 50% for others; PTRG under review for flexi-load
   % [Q10c] say "a compressed week may fit where WFH does not" (fixer pass 7 Oct 2026, employer-11; fact_audit s3 item 6)
   \textbf{Firms}, with SNEF: review which tasks need presence; a compressed week may fit where WFH does not; use flexi-load where fewer hours, not location, address the need. Job Redesign+: up to 70\% of qualifying costs for SMEs, 50\% for others. \textbf{MOM:} PTRG review.}
  {% [P11] EWTP from Mar 2026, linked to FWAs at COS 2026
   EWTP (Mar 2026), linked by MOM to FWAs.}
  {Not one-size-fits-all; not a quota.}
  {% [B1] fairness top challenge at 200 or fewer staff, 28%; [P5] COS 2026 speech: under half who needed part-time got it, 2024
   Fairness to staff without FWAs (top challenge, $\le$200 staff); part-time access (COS 2026: under half who needed it got it).}
  {\SourceNote{Associations; asked after the choices. \hyperlink{B-flexiload}{\textcolor{Navy}{Flexi-load}}.}%
   \BackToTwo{M-options}{Options}{M-where}{Reservations}}

% [S7] 3-4 WFH days is a tested attribute level; W5 M-options row C
\OptionFrame{B-optC}{Option C: a 3--4-day WFH pilot where roles allow}
  {\textbf{Zone:} C, far apart\quad\textbf{Strength:} exploratory}
  {
   Employers would require pay cuts of about 36\% for 3--4 WFH days and 49\% for fully remote work (PMEs would give up 10\% and 12\%). Fully remote draws about 61\% where firms worry that creativity suffers apart, against 36\% elsewhere (a weak association, p = .14; preliminary). Abroad, a six-month hybrid trial cut quits by a third.}
  {% W5 M-options owner: volunteer firms; independent evaluation; Lead decision 10: roles without presence needs
   \textbf{Volunteer firms}, with SNEF and NTUC: a 3--4-day WFH pilot in roles without presence needs; outcomes agreed upfront; independent evaluation; stop criteria set in advance. Convenor, funding and publication rule: for the partners.}
  {% [V6] managers' perceived productivity effect -2.6% to +1.0%
   Trials abroad (managers' productivity view: $-$2.6\% to +1.0\%).}
  {% [L1] fully remote evidence mixed (FACTCHECK F25)
   Not a push for fully remote work (mixed evidence; some studies find lower productivity); not a mandate; no pay change.}
  {Output, retention, collaboration; whether managers' concerns are borne out.}
  {\SourceNote{Bloom, Han \& Liang (2024); Barrero, Bloom \& Davis (2023); Emanuel \& Harrington (2024).}%
   \BackToTwo{M-options}{Options}{M-meet}{Two sides}}

\OptionFrame{B-optD}{Option D: why-not-request and partial-approval items}
  {\textbf{Zone:} all (measurement)\quad\textbf{Strength:} gap documented}
  {% [P2] no administrative counts; [P3] 18.5% made a formal request, data for 2024; [D9] misread item; [B6] no why-not item
   % [D7] 4.4 [7.2] declined, 32 [42] partly; [D7b] lower bound; [M23] 44.6 [37.1]; [M27] more formal requests, not fewer
   % fixer pass 7 Oct 2026 (mom_ps-02 b): [P3] MOM publishes request rates; [B6] our survey had no why-not item; whether MRSD's
   % questionnaire already asks why is a PI gate before the meeting (REPORT.md)
   Firms need not report requests to MOM; in MOM's survey, 18.5\% of those needing an FWA made a formal request (data for 2024). MOM publishes request rates, not reasons; our survey had no why-not item, and our request item was partly misread. Of PMEs reporting a request, about 4\% were declined (a lower bound) and 32\% partly approved. 45\% expect a career penalty; it goes with more requests, not fewer.}
  {% W5 M-options owner: MRSD checks the items against its survey, with NTUC and SNEF (fixer pass 7 Oct 2026, mom_ps-02);
   % COMMENTS slide 20: first deliverable restricted to the two items
   \textbf{MOM (MRSD)}, with NTUC and SNEF: check its survey, then add the items not yet asked: why those who need an FWA do not ask, and whether a request was approved in full, with changes, or declined (\hyperlink{B-optD-items}{\textcolor{Navy}{draft}}).}
  {MOM's employee survey; partners' outreach channels.}
  {Not a new reporting duty for firms; not an evaluation of the Guidelines.}
  {% [D6] junior 29 [25] familiar; [P15] about 6% of FWA users reported a setback in 2025, experience not expectation
   Partial approvals; junior PMEs' awareness (29\% familiar); expected penalty vs MOM's reported setbacks (\ap 6\% of FWA users, 2025).}
  {\SourceNote{Surveyed samples, unweighted.}%
   \BackToTwo{M-options}{Options}{M-barriers}{Barriers}}

\begin{frame}[t,label=B-optD-items]{Option D: draft survey items}
  \FrameTop
  \setbeamerfont{itemize/enumerate body}{size=\fontsize{9.5}{11.6}\selectfont}
  % COMMENTS slide 20 and s4: the first deliverable is the why-not-request and partial-approval items,
  % with an automatic-policy category; [D9] misread request item; [B6] no why-not item; [P2] no administrative counts
  \begin{columns}[T,onlytextwidth]
    \begin{column}{0.48\textwidth}
      \PanelTitle{Employee}{First step (employees)}
      \begin{enumerate}\setlength{\itemsep}{0.08cm}
        \item \textbf{Why not ask:} if you needed an FWA but did not request one, why not? (career concern; did not know how; expected refusal; automatic)
        \item \textbf{Outcome:} approved as requested / with changes / declined.
      \end{enumerate}
      \vspace{0.12cm}
      \PanelTitle{Employee}{Also (employees)}
      \begin{enumerate}\setcounter{enumi}{2}\setlength{\itemsep}{0.08cm}
        \item Formal request: add ``applies automatically''.
        \item Guidelines awareness; would flexible work affect your career?
      \end{enumerate}
    \end{column}
    \begin{column}{0.48\textwidth}
      % fixer pass 7 Oct 2026 (employer-04): establishment items only if the partners agree
      \PanelTitle{Employer}{Later, if partners agree (establishments)}
      \begin{enumerate}\setcounter{enumi}{4}\setlength{\itemsep}{0.08cm}
        \item Any formal FWA request in the past 12 months?
        \item Share approved as requested, with changes, or declined?
        \item Roles needing on-site presence: which arrangements are offered?
      \end{enumerate}
    \end{column}
  \end{columns}
  \SourceNote{Draft for MRSD, to check against EWC and CoE questionnaires; item 1 appears new. Employee items first; any establishment items would sit in existing MOM surveys, for the partners to decide. \hyperlink{B-optD}{\textcolor{Navy}{Option D}}.}
  \BackToTwo{M-options}{Options}{M-conclusion}{Conclusion}
\end{frame}

% ==========================================================================
% Technical appendix
% ==========================================================================
\begin{frame}[t,label=B-why-org-tests]{Technical: all 11 employer split tests}
  \FrameTop % first body token (centring; never a frame subtitle)
  % Appended sample (n = 397): ratio tests of equal location WTPs across groups (group-specific pay slopes), data/wtp_deck.csv (ratio_p);
  % q = Benjamini-Hochberg over the 11 tests; 3-group splits use the joint 6-df test
  {\TableFont\setlength{\tabcolsep}{4pt}\renewcommand{\arraystretch}{1.0}%
  \begin{tabularx}{\linewidth}{@{}>{\RaggedRight}X>{\centering}p{3.2cm}>{\centering}p{1.6cm}>{\centering\arraybackslash}p{1.6cm}@{}}
    \toprule
    \textbf{Split} & \textbf{Firms} & \textbf{p} & \textbf{q}\\
    \midrule
    No WFH now & 96 / 301 & .023 & .14\\
    Presence need & 137 / 260 & .025 & .14\\
    Creativity concern & 201 / 196 & .14 & .43\\
    Concern index & 260 / 137 & .19 & .43\\
    Written policy & 276 / 121 & .19 & .43\\
    Hiring role (3 groups) & 84 / 122 / 170 & .36 & .61\\
    Productivity concern & 223 / 174 & .39 & .61\\
    Coordination concern & 213 / 184 & .47 & .62\\
    Salary offset & 187 / 210 & .51 & .62\\
    Firm size (3 bands) & 175 / 140 / 82 & .63 & .66\\
    Tech readiness (3 levels) & 29 / 209 / 159 & .66 & .66\\
    \bottomrule
  \end{tabularx}\par}
  % EB s9 l.182: presence and concern items sit in Block 5, after the DCE block
  \SourceNote{Technical appendix, not for projection. p: test that the three location pay-cut values (1--2 days, 3--4 days, fully remote) are equal across the groups, with a separate pay slope per group; q: Benjamini--Hochberg over the 11 tests. Items asked after the choices: preliminary associations.}
  \BackTo{M-where}{Reservations}
\end{frame}

% Round 5: in the technical appendix (presenter preparation, not for projection); wording follows the
% round-5 language rules (MANDATE, Evidence and language rules; ROUND5_SPEC s6).
\begin{frame}[t,label=B-language]{Public wording}
  \FrameTop
  \setbeamerfont{itemize/enumerate body}{size=\fontsize{9.5}{11.4}\selectfont}
  \PanelTitle{Navy}{We say}
  \begin{itemize}\setlength{\itemsep}{0.02cm}
    % [E1, E2] about 7-10, indicative; never a pay cut
    \item PMEs strongly prefer hybrid work: about 7--10\% of pay, indicative, never a pay cut.
    % [Q1, Q2, Q5] about six in ten; EB s3b plain-English reading, choices not people
    \item Employers would require larger pay cuts than PMEs would give up for most arrangements; for the compressed week the two are about equal.
    % [R5, Q5] no detectable penalty; [R10, R12] fully remote ranked last (both hold in the appended sample)
    \item Hiring: no detectable gap for a compressed week; fully remote draws the largest cut (a ranking).
    % [M10] presence and creativity contrasts; EB s9: free-slope check pending
    \item Reservations: 1--2 WFH days where most roles need presence; remote with the creativity concern; preliminary (asked after the choices).
    % Lead decision 2; mandate: compressed-week conclusion wording
    \item Compressed week: most promising common-ground test; evidence moderate; operating costs unmeasured.
    % [M23] more than a third; [D7, D7b] lower bound
    \item Over a third of surveyed PMEs expect a career penalty; refusals are a lower bound.
  \end{itemize}
  \vspace{0.06cm}
  \PanelTitle{Navy}{We avoid}
  \Body{Pay values as a proposed pay cut, or employer cut sizes as a price (read the ranking); a request rate; ``easy'', ``low cost'', ``no cost'', ``indifferent''; associations or nulls as causes or proof; Guidelines effects on provision, requests or approvals.}
  \BackToTwo{M-meet}{Two sides}{M-conclusion}{Conclusion}
\end{frame}

% Round 5 (ROUND5_SPEC s5; COMMENTS B-org-pct): model-dependent header; never linked from a short-path
% main slide (W3 links B-org-salary from M-employers instead). [R1-R9] weakly identified; EB s3 header; EB s8
\begin{frame}[t,label=B-org-pct]{Employer \% of salary}
  \FrameTop % first body token (centring; never a frame subtitle)
  {\centering\Caveat{Model-dependent; not a pay recommendation and not a main finding.}}
  \Tab{T_org_pct}
  \SourceNote{Firmer result: the ranking of the pay cuts.}
  \BackTo{M-employers}{Employers}
\end{frame}

% Round 5 (ROUND5_SPEC s5; COMMENTS B-org-pct and B-gap): model-dependent header; never linked from a
% short-path main slide. [G1-G6] gap, indicative; [E3] the T_gap remote row is vs a long commute, flagged in its own footnote
\begin{frame}[t,label=B-gap]{Gap, \% of pay: indicative}
  \FrameTop % first body token (centring; never a frame subtitle)
  {\centering\Caveat{Model-dependent; not a pay recommendation and not a main finding.}}
  \Tab{T_gap}
  % EB s3 header: employer magnitudes weakly identified; Lead decision 3: no cross-side ruler
  \BackTo{M-meet}{Two sides}
\end{frame}

% ==========================================================================
% PME groups
% ==========================================================================
% [E18] gender p = .43
\BackupFig{B-ra-gender}{Estimates by gender}{F_ra_gender}{5.4cm}{M-groups}{PME groups}{No statistically significant gender difference (joint test p = .43). \hyperlink{B-emp-groups}{\textcolor{Navy}{All groups}}.}

% [E18] age p = .09
\BackupFig{B-ra-age}{Estimates by age}{F_ra_age}{5.4cm}{M-groups}{PME groups}{Ages 19--39 vs 40+; not statistically significant (joint test p = .09). \hyperlink{B-emp-groups}{\textcolor{Navy}{All groups}}.}

% [E18] caregiving p = .66
\BackupFig{B-ra-care}{Estimates by caregiving}{F_ra_care}{5.4cm}{M-groups}{PME groups}{No caregiving vs any; not statistically significant (p = .66). \hyperlink{B-emp-groups}{\textcolor{Navy}{All groups}}.}

% [E18] role p = .15
\BackupFig{B-ra-role}{Estimates by role level}{F_ra_role}{5.4cm}{M-groups}{PME groups}{Junior, middle and senior; not statistically significant (p = .15). \hyperlink{B-emp-groups}{\textcolor{Navy}{All groups}}.}

% [E18b] education p = .56
\BackupFig{B-ra-edu}{Estimates by education}{F_ra_edu}{5.4cm}{M-groups}{PME groups}{Three education levels; not statistically significant (p = .56). \hyperlink{B-emp-groups}{\textcolor{Navy}{All groups}}.}

% ==========================================================================
% PME checks
% ==========================================================================
% [E15] 73 same-side respondents; within about 1 point
\BackupFig{B-ra-sameside}{Without same-side respondents}{F_ra_sameside}{5.4cm}{M-robust}{Robustness}{Drops 73 PMEs who always picked one side; location and compressed-week values move $\le$1 point. \hyperlink{B-emp-robust}{\textcolor{Navy}{PME checks}}.}

% [E15] trimming 1%, n = 1,007
\BackupFig{B-ra-trim}{Trimming extreme salaries}{F_ra_trim}{5.4cm}{M-robust}{Robustness}{Drops the top and bottom 1\% of salaries. \hyperlink{B-emp-robust}{\textcolor{Navy}{PME checks}}.}

\end{document}
